% Source: dropout_as_regularization.tex
\begin{notechapter}
\chapter{Dropout as a Regularization Method}\label{ch:dropout-regularization}

\section{Introduction}
Dropout is a regularization technique commonly used in neural 
networks to prevent overfitting. This document explains how 
dropout works, why it is effective, and how it connects to 
Bayesian model averaging.

\section{Concept of Dropout}
During training, dropout randomly "drops out" (i.e., sets to 
zero) a fraction of the neurons in a layer with a specified 
probability \( p \) (typically 0.5 for hidden layers). This 
means each neuron has a probability \( p \) of being 
temporarily removed from the network during each training 
iteration.

\section{Regularization Effect}
\begin{itemize}
    \item \textbf{Preventing Co-Adaptation:} Dropout prevents 
    neurons from becoming overly reliant on specific other 
    neurons. By randomly dropping neurons, the network is 
    forced to learn robust features that are not dependent on 
    any single neuron or small group of neurons.
    
    \item \textbf{Ensemble Effect:} Dropout can be interpreted 
    as training an ensemble of many different sub-networks 
    (each with a random subset of neurons) simultaneously. 
    During inference, the full network approximates averaging 
    the predictions of all these sub-networks.
\end{itemize}

\section{Mathematical Interpretation}
During training, the output of a neuron is scaled by 
\(\frac{1}{1-p}\) to compensate for the dropped neurons. This 
ensures the expected value of the neuron's output remains the 
same during training and inference. For example, if \( p = 0.5 
\), the output is scaled by 2.

\section{Why Scaling is Necessary}
Without scaling, the expected output of a neuron during 
training would be smaller than during inference. Scaling by 
\(\frac{1}{1-p}\) ensures consistency. Mathematically:
\[
E[\text{output}] = p \cdot 0 + (1-p) \cdot \left( \frac{x}{1-p} 
\right) = x
\]
This matches the neuron's output during inference.

\section{Dropout During Training vs. Inference}
\begin{itemize}
    \item \textbf{Training:} Dropout is applied to randomly 
    drop neurons, and outputs are scaled by \(\frac{1}{1-p}\).
    
    \item \textbf{Inference:} Dropout is turned off, and the 
    full network is used. No scaling is applied.
\end{itemize}

\section{Dropout as Approximate Bayesian Inference}
Gal and Ghahramani (2016) showed that training a network with dropout
can be interpreted as approximate \emph{variational} Bayesian inference:
the random masks define a (very particular) approximate posterior
distribution over the weights, and the dropout objective is, up to
constants, the corresponding variational objective. It is an
approximation: dropout does not sample from the true posterior.

\subsection{Bayesian Model Averaging}
Under this interpretation, predictions should average over masks,
which corresponds to (approximate) Bayesian model averaging:
\[
\text{Inference Output} \approx \mathbb{E}_{m \sim p(m)}[f(x; 
\theta, m)]
\]
where \( f(x; \theta, m) \) is the output of the network with 
parameters \( \theta \) and dropout mask \( m \).

\subsection{Why the Full Network Approximates the Expectation}
The scaling during training ensures that the expected value
of each neuron's input is preserved. Using the full network at
inference (the ``weight scaling inference rule'') replaces the
average of the outputs by the output of the average network:
\[
\mathbb{E}_{m \sim p(m)}[f(x; \theta, m)] \approx f(x; \theta, 
\text{no dropout})
\]
This is exact for a linear network, but only an approximation once
non-linearities are involved, since in general
\( \mathbb{E}[\phi(Z)] \neq \phi(\mathbb{E}[Z]) \). \emph{Monte Carlo dropout}
keeps dropout active at test time and averages several stochastic
forward passes, which estimates the expectation directly and also
provides an estimate of the predictive uncertainty.

\section{Practical Implementation}
In deep learning frameworks like PyTorch, dropout is 
implemented as a layer that behaves differently during 
training and inference. For example:
\begin{verbatim}
dropout_layer = torch.nn.Dropout(p=0.5)  # 50% dropout
output = dropout_layer(input)  # Training: applies dropout
dropout_layer.eval()  # Switch to evaluation mode
output = dropout_layer(input)  # Inference: no dropout
\end{verbatim}

\section{Conclusion}
Dropout is a powerful regularization technique used only 
during training to prevent overfitting. It introduces 
randomness, prevents co-adaptation, and approximates an 
ensemble of sub-networks. During inference, dropout is 
disabled to ensure stable and deterministic predictions. 
Additionally, dropout can be interpreted as approximate
variational Bayesian inference; the deterministic full network
is a cheap approximation of the corresponding model average.
\end{notechapter}
